
The source and destination buffers for all data movement operations and the
symmetric data objects for all synchronization operations in this section shall
reside in the \ac{GPU} symmetric heap. Using local \ac{GPU} memory, local CPU
memory, or memory from the default symmetric heap as source, destination, or
synchronization operands for GPU-centric communication operations results in
\textbf{undefined behavior}.

For routines with a \VAR{pe} argument, the \VAR{pe} argument shall identify a
\ac{PE} in the team associated with the given \VAR{ctx} when provided, or the
default context otherwise. Passing an invalid \ac{PE} number results in
\textbf{undefined behavior}.

Each communication routine is provided at three thread granularities:

\begin{longtable}{|p{3cm}|p{3cm}|p{8.5cm}|}
\hline
\textbf{Suffix} & \textbf{Granularity} & \textbf{Description} \\
\hline
\endfirsthead
\hline
\textbf{Suffix} & \textbf{Granularity} & \textbf{Description} \\
\hline
\endhead
\emph{(none)} & Thread & A single \ac{GPU} thread initiates the operation. \\
\hline
\texttt{\_tg} & Thread group & All threads in a thread group collectively initiate
the operation (e.g., a thread block or workgroup). \\
\hline
\texttt{\_htg} & Hardware thread group & All threads in a hardware thread group
collectively initiate the operation (e.g., a warp or wavefront). \\
\hline
\end{longtable}

The thread group (\texttt{\_tg}) and hardware thread group (\texttt{\_htg})
variants are collective across the threads in the specified group. All
participating threads shall call the routine with the same arguments. If any
thread in the group does not participate, or if arguments differ across threads,
the behavior is undefined.

The atomic memory operations (\cref{subsec:gc_amo}) and point-to-point
synchronization routines (\cref{subsec:gc_p2p}) are provided as typed variants,
following the naming convention of the \openshmem 1.6 specification. In the name
of a typed routine, \TYPENAME{} is the type-name component that identifies the
operand type, while \TYPE{} is the corresponding C type used in the routine
signature. The supported types and the \TYPE{}-to-\TYPENAME{} correspondence are
those defined by the \openshmem 1.6 specification for the relevant operation
category. For example, the type \TYPE{} \KEYWORD{unsigned long long} has
\TYPENAME{} \texttt{ulonglong}, giving rise to routines such as
\FUNC{shmemg\_ulonglong\_atomic\_add}.

